% Part: computability
% Chapter: recursive-functions
% Section: composition

\documentclass[../../../include/open-logic-section]{subfiles}

\begin{document}

\olfileid{cmp}{rec}{com}
\olsection{التركيب}

إذا كانت $f$ و$g$ دالتين أحاديتي الحجة في الأعداد الطبيعية، أمكننا
تركيبهما: $h(x)=f(g(x))$. وعندئذ تعرف الدالة الجديدة~$h(x)$
\emph{بالتركيب} من الدالتين $f$ و~$g$. ونريد تعميم ذلك على الدوال ذات
أكثر من حجة.

إليك إحدى الطرق: افترض أن $f$ دالة ذات $k$ حجج، وأن $g_0$، …،
$g_{k-1}$ دوال عددها $k$ وكلها ذات $n$ حجة. يمكننا عندئذ تعريف دالة
جديدة~$h$ ذات $n$ حجة على النحو الآتي:
\[
h(x_0, \dots, x_{n-1}) =
f(g_0(x_0, \dots, x_{n-1}), \dots, g_{k-1}(x_0, \dots, x_{n-1}))
\]
إذا كانت $f$ وجميع $g_i$ قابلة للحساب، كانت $h$ كذلك. فلحساب
$h(x_0,\dots,x_{n-1})$، نحسب أولًا القيم
$y_i=g_i(x_0,\dots,x_{n-1})$ لكل $i=0$، …،~$k-1$. ثم ندخل هذه
القيم في $f$ فنحسب
$h(x_0,\dots,x_{n-1})=f(y_0,\dots,y_{k-1})$.

قد يبدو هذا توصيفًا شديد التقييد لما يحدث حين نحسب دالة جديدة باستعمال
دوال موجودة. فمن جهة، لا نستعمل أحيانًا جميع حجج دالة ما، كما حين عرفنا
$g(x,y,z)=\Succ(z)$ لاستعمالها في التعريف بالعودية البدائية لـ~$\Add$.
افترض أنه يسمح لنا باستعمال الدوال الآتية:
\[
\Proj{n}{i}(x_0, \dots, x_{n-1}) = x_i
\]
تسمى الدوال~$\Proj{n}{i}$ \emph{دوال الإسقاط}؛ فـ$\Proj{n}{i}$ دالة
ذات $n$ حجة. وعندئذ يمكن تعريف $g$ بـ
\[
g(x, y, z) = \Succ(\Proj{3}{2}(x, y, z)).
\]
هنا تؤدي الدالة الأحادية الحجة $\Succ$ دور $f$، ولذلك $k=1$. ولدينا
دالة واحدة ثلاثية الحجج $\Proj{3}{2}$ تؤدي دور $g_0$. والنتيجة دالة
ثلاثية الحجج تعيد لاحق الحجة الثالثة.

وتتيح لنا دوال الإسقاط أيضًا تعريف دوال جديدة بإعادة ترتيب الحجج أو
توحيد بعضها. فمثلًا يمكن تعريف الدالة $h(x)=\Add(x,x)$ بـ
\[
h(x_0) = \Add(\Proj{1}{0}(x_0),\Proj{1}{0}(x_0)).
\]
هنا $k=2$ و$n=1$، وتؤدي $\Add$ دور $f(y_0,y_1)$، ويؤدي
$\Proj{1}{0}(x_0)$، وهو دالة الإسقاط الأحادية الحجة (أي دالة الهوية)،
دوري $g_0(x_0)$ و$g_1(x_0)$ معًا.

وإذا كانت $f(y_0,y_1)$ دالة لدينا من قبل، أمكننا تعريف الدالة
$h(x_0,x_1)=f(x_1,x_0)$ بـ
\[
h(x_0, x_1) = f(\Proj{2}{1}(x_0, x_1),\Proj{2}{0}(x_0, x_1)).
\]
هنا $k=2$ و$n=2$، وتؤدي $\Proj{2}{1}$ و~$\Proj{2}{0}$ دوري
$g_0$ و$g_1$ على الترتيب.

وقد يقلقك أيضًا اشتراط أن تكون $g_0$، …،~$g_{k-1}$ كلها ذات عدد
الحجج نفسه~$n$. (تذكر أن \emph{رتبة} الدالة هي عدد حججها؛ فالدالة ذات
$n$ حجة رتبتها~$n$.) غير أن إضافة دوال الإسقاط توفر المرونة المطلوبة.
فمثلًا، افترض أن $f$ و~$g$ دالتان ثلاثيتا الحجج وأن $h$ دالة ثنائية
الحجة معرفة بـ
\[
h(x,y) = f(x,g(x,x,y),y).
\]
يمكن إعادة كتابة تعريف~$h$ بدوال الإسقاط هكذا:
\[
h(x,y) = f(\Proj{2}{0}(x,y), g(\Proj{2}{0}(x,y), \Proj{2}{0}(x,y),
\Proj{2}{1}(x,y)), \Proj{2}{1}(x,y)).
\]
وعندئذ تكون $h$ تركيب $f$ مع $\Proj{2}{0}$ و$l$ و
~$\Proj{2}{1}$، حيث
\[
l(x,y) = g(\Proj{2}{0}(x,y),\Proj{2}{0}(x,y),\Proj{2}{1}(x,y)),
\]
أي إن $l$ تركيب $g$ مع $\Proj{2}{0}$ و$\Proj{2}{0}$ و
~$\Proj{2}{1}$.

\end{document}
